Temperatur gehört zu den wichtigsten Prozessgrößen der industriellen Messtechnik. Thermoelemente eignen sich insbesondere für robuste Messungen über große Temperaturbereiche, liefern jedoch nur kleine Thermospannungen und benötigen deshalb eine geeignete Signalaufbereitung. Ziel des Versuchs ist es, die gesamte Messkette vom Sensor bis zur digitalen Anzeige experimentell zu untersuchen. Dazu werden die Signalverstärkung und die ADC-Ausnutzung dimensioniert, die Messkette gegen ein Referenzsystem kalibriert und die Empfindlichkeit gegenüber hochfrequenten Störungen analysiert. Abschließend wird ein passiver RC-Tiefpass entworfen und hinsichtlich Dämpfung und Dynamik bewertet.


\begin{figure}[htbp]
    \centering
    \includegraphics[width=\linewidth]{V2_1_char_curve.pdf}
\end{figure}